% =====================================================================
%  Reproducibilidad de la corrida canónica
%  Anexos
%
%  ENCARGO: Comandos exactos, entorno congelado, huellas de los artefactos y compuertas de verificación.
%
%  Reglas del documento:
%   - Una idea = una subseccion = una figura.
%   - Toda figura declara su procedencia con \fuente{...}.
%   - Toda cifra sale del canon de agosto; ver GUIA.md.
%   - M1 y M3 se reportan siempre en paralelo.
%   - Ningun superlativo sin releer el CSV de la frontera correcta.
% =====================================================================
\section{Reproducibilidad de la corrida canónica}
\label{anx:repro}
\justifying

\subsection{Cómo se localizan y se leen los archivos crudos}
\label{sub:repro-lectura}

El Capítulo~\ref{sec:prep} resume en una sola etapa lo que el cargador
hace para llegar a la serie de cada equipo. Se llama así aquí a la pieza
que recorre los archivos, los lee y entrega la serie horaria de un equipo. El detalle queda aquí, porque
describe el programa y no el dato. De las decisiones que siguen, casi
ninguna llega a ejercerse sobre este conjunto, y la única que sí se
ejerce no cambia ningún valor.

\textbf{La búsqueda de carpetas tolera diferencias de escritura.} Los
nombres con que llegan las carpetas no son uniformes, y el de los
medidores de la Universidad CESMAG viene con una errata de origen, escrita
\code{eletricMeter}. Se conserva tal cual, para reproducir la estructura de
la fuente en vez de corregirla y perder la trazabilidad. La tolerancia, sin
embargo, no llega a hacer falta: las diez carpetas que el modelo emplea
coinciden exactamente con lo declarado en la configuración.

\textbf{Cada archivo se decodifica probando cuatro codificaciones en
orden}: \code{utf-8-sig}, \code{utf-8}, \code{cp1252} y \code{latin-1}. El
orden importa porque \code{latin-1} acepta cualquier secuencia de bytes y
nunca falla, de modo que probarla antes que \code{cp1252} la haría ganar
siempre y corrompería los caracteres del rango \code{0x80}--\code{0x9F},
que es donde viven las tildes.

Tampoco esa cascada se ejerce. Los 73 archivos del árbol resuelven en el
primer candidato, y ninguno contiene un solo byte por encima de 127. Se
documenta porque el orden es una decisión deliberada del cargador y porque
un conjunto de datos posterior sí podría ejercerla, no porque haya
cambiado una cifra de este trabajo.

\textbf{La fecha y el valor se convierten en modo tolerante.} Lo que no se
deja interpretar queda nulo en vez de detener la carga, y las dos columnas
no reciben el mismo trato: una lectura sin fecha válida se descarta entera,
porque no tiene dónde colocarse, mientras que una con fecha válida y valor
ilegible se conserva en blanco para que la limpieza la rellene más
adelante. Ninguna de las dos conversiones llega a descartar nada, porque
sobre los \num{2813305} registros de los diez medidores que el modelo
emplea no hay una sola fecha ilegible ni un solo valor no numérico.

\textbf{Los instantes repetidos se funden promediando.} La fusión ocurre
dentro de cada archivo, antes de alinear los tramos de un mismo medidor, y
su motivo es de forma y no de contenido: el índice ha de ser único para que
los tramos se puedan superponer sobre un mismo eje de tiempo. La
Tabla~\ref{tab:repro-dup} recoge dónde ocurre.

\begin{table}[H]
 \centering
 \caption{Instantes con más de una lectura, por institución y frontera. Las
 instituciones que no aparecen no tienen ninguno.}
 \label{tab:repro-dup}
 \begin{tabular}{@{}l l r r@{}}
 \toprule
 \textbf{Institución} & \textbf{Frontera} &
 \textbf{Instantes repetidos} & \textbf{Máximo en un instante} \\
 \midrule
 UCC & M1 & \num{3071} & 11 \\
 UCC & M3 & \num{3551} & 73 \\
 \addlinespace
 Mariana & M1 & \num{28} & 2 \\
 Mariana & M3 & \num{71} & 4 \\
 \bottomrule
 \end{tabular}
\end{table}

Las tres reglas de fusión posibles, promediar, quedarse con la primera o
quedarse con la última, dan el mismo resultado sobre este conjunto, porque
las lecturas que comparten instante coinciden hasta el último dígito en la
totalidad de los casos. Lo que sí informa es la concentración: la UCC
aporta el \pct{98.5} de los instantes repetidos, y su medidor secundario
llega a volcar 73 lecturas idénticas en un mismo instante. Es un rasgo del
equipo y no del método, y no altera ninguna cifra.

\textbf{La unidad se convierte al final, y ahí el orden solo importa para
el recorte.} El cargador promedia a paso horario y solo después divide
entre mil la serie del inversor. Para la división el orden es indiferente,
porque dividir entre una constante y promediar se pueden intercambiar sin
que el resultado cambie; comprobado sobre los siete inversores, invertirlo
produce diferencias a partir de la decimoquinta cifra decimal, que es ruido
de redondeo. Con el recorte a valores no negativos no ocurre lo mismo:
recortar y promediar no se pueden intercambiar, porque una lectura negativa
dentro de una hora de media positiva sobreviviría al promedio y rebajaría
el valor horario. El cargador recorta sobre la serie ya horaria, de modo
que ese caso podría darse; sobre este conjunto no llega a existir, porque
ninguna de las \num{793981} lecturas crudas de los siete inversores es
negativa.

\textbf{Los tres tipos de medidor llevan otro nombre en el canon.} El
Capítulo~\ref{sec:prep} los llama neto, neto parcial y bruto, que es el
vocabulario del texto y de las figuras. En la configuración del cargador y
en las tablas de datos que acompañan a cada figura figuran como
\code{net}, \code{net\_partial} y \code{gross}. Se anota la
correspondencia aquí para que quien audite el código no la tenga que
adivinar.

\textbf{El orden exacto} en que se aplica todo lo anterior es el siguiente:
localizar la carpeta del equipo, decodificar cada archivo, convertir la
fecha y el valor, fundir los instantes repetidos, alinear los tramos del
medidor sobre un mismo eje de tiempo, recortar al horizonte, promediar a
paso horario y, en último lugar, convertir la unidad.

\textbf{Cómo se lee la figura del recorrido.} El
Capítulo~\ref{sec:prep} traza una hora real desde el archivo hasta la serie
horaria, y conviene decir qué hay en cada banda de la
Figura~\ref{fig:prep-archivo-serie}. Arriba, los archivos del medidor y del
inversor sobre el calendario, con el horizonte de estudio sombreado. En el
centro, las lecturas a paso de dos minutos de esa hora: el medidor entrega
38 filas que al depurar quedan en 30 instantes, ocho de los cuales traen la
lectura duplicada, y el inversor entrega 30 lecturas en vatios, ninguna
negativa, de modo que el recorte a cero no llega a actuar. Abajo, la media
de cada panel ocupa su lugar entre las 24 horas del día.

Fundir los duplicados quita una copia y no rellena nada. Los instantes que
quedan caen en los treinta tramos de dos minutos de la hora sin faltar
ninguno, de modo que el ejemplo ilustra el caso en que la medida cubre el
intervalo entero y no hay tramo que suponer.

\subsection{El criterio del regulador para la lectura que falta}
\label{sub:repro-norma}

El Capítulo~\ref{sec:prep} invoca la norma como criterio y no como
obligación. Aquí queda el articulado que lo sostiene, porque la analogía
solo vale si se puede comprobar.

El artículo 38 de la Resolución CREG 038 de 2014 se activa mientras se
reparan o reponen los elementos de un sistema de medición en falla o
hurtado, y ordena estimar la lectura por dos vías distintas según la
frontera. Si la frontera reporta al administrador del sistema de
intercambios comerciales, integrando la medida de potencia activa cuando
esa medida está bajo la supervisión de un centro de control, o recurriendo
a curvas típicas. Si no reporta, conforme a la Ley 142 de 1994 y al
artículo 31 de la Resolución CREG 108 de 1997, que ordena establecer el
consumo con base en los consumos promedios de otros períodos del mismo
usuario, en los de usuarios en circunstancias similares o en aforos
individuales.

Dos precisiones importan. La primera es que \textbf{ninguno de los medios
que el artículo enumera consiste en asignar cero}, y conviene decirlo así
porque es una afirmación sobre la lista y no sobre una prohibición que la
norma no formula. La segunda es que la vía aplicable por analogía a este
trabajo es la segunda y no la primera: integrar la potencia activa significa
allí leer otro canal que sí funciona, mientras que estimar por promedios de
otros períodos del mismo usuario es exactamente lo que aquí se hace con la
hora incompleta.

Nada de esto obliga. Los equipos del proyecto son instrumentación de
investigación y no fronteras comerciales, y lo que interrumpe el registro no
es un sistema de medición en falla sino un corte de telemetría.

El texto del artículo 38 se verificó por dos vías independientes que
coinciden palabra por palabra. Se deja anotado porque la publicación oficial
en línea de la resolución se trunca dentro del artículo 28 y no alcanza al
38, de modo que quien quiera comprobarlo no lo encontrará en la fuente
evidente.

\subsection{Cómo se lee la figura de la profundidad de la lectura}
\label{sub:repro-profundidad}

La Figura~\ref{fig:prep-gradacion} del Capítulo~\ref{sec:prep} resume cada
uno de los cinco medidores de la frontera principal en un renglón, y la
distancia entre sus dos marcas es la gradación que la subsección describe.
Fila por fila:

En Udenar el medidor baja bastante más abajo del cero de lo que marca de
ordinario, es decir, la inversión del flujo no es allí un episodio sino el
estado normal de las horas de sol.

En Mariana y en la UCC el descenso no llega siquiera a igualar la lectura
mediana de su propio circuito, de manera que el flujo se invierte solo
cuando la producción culmina y la carga afloja.

En el HUDN el recorrido entero se queda por encima de \uni{6}{kW}, y en
CESMAG termina en \uni{0.195}{kW}, es decir, roza el cero sin cruzarlo,
porque detrás de esos dos medidores no genera nada.

\placeholderfig{Anexo pendiente de redacción. Comandos exactos, entorno congelado, huellas de los artefactos y compuertas de verificación.}
